% ============================================================
% VIII. CONCLUSION
% ============================================================
\section{Conclusion}
\label{sec:conclusion}

We have presented URDFApproxGeom, an integrated open-source toolchain for automatic conversion of mesh-based URDF collision geometry into convex, sphere-tree, and capsule approximations.
The toolchain provides a unified CLI and Python API, named presets, validation metrics with pass/fail gates, and qualitative visualization support, making it immediately usable in existing robotics workflows.

The main methodological contribution is the automatic capsule-output pipeline, which produces low-count analytic capsule primitives whose vertex coverage, tightness, and runtime can be measured and tuned through the accompanying validation workflow. The pipeline proceeds through principal-axis sectioning, adaptive circle fitting, collinear merging, vertex-coverage-aware growth, endpoint tightening, inflation-based splitting, and nested-capsule pruning.
Each capsule is emitted as a URDF-compatible cylinder-plus-two-sphere decomposition alongside a JSON sidecar with the canonical analytic parameters.

\subsection*{Future Work}

Several directions for future development are evident:

\paragraph{Broader robot benchmarks.}
The current evaluation focuses on the Franka FR3 arm.
Applying the pipeline to a wider range of robot morphologies---including mobile manipulators, humanoids, parallel kinematics, and non-tubular links---would reveal generalizability limits and guide algorithmic improvements.

\paragraph{Topology-aware capsule decomposition.}
The PCA-based sectioning assumes a single dominant axis per link.
For branched or non-elongated geometry, algorithms that discover multiple local axes or use volumetric over-segmentation before per-segment capsule fitting could improve approximation quality.

\paragraph{URDF format extension.}
Advocating for and implementing a native capsule element in the URDF standard would eliminate the need for cylinder-sphere decomposition, enabling simulators and planners to directly use the capsule distance function.

\paragraph{Tighter integration with collision libraries.}
The JSON sidecar format is a step toward making primitive parameters reusable, but tighter integration with collision-checking libraries (e.g., FCL~\cite{pan2012fcl} and its derivatives, or the Drake geometry framework) would allow the generated capsules to be used natively without manual sidecar parsing.

\paragraph{Learning-based fitting.}
While the current pipeline is entirely geometric, a learned approach could potentially predict capsule configurations from mesh topology, though generalizability beyond the training distribution and the cost of generating labeled training data remain open challenges.
